\chapter{Tools Adopted}
\label{ch:adopted}

\section{GENESYS Client API}
\label{sec:genesys-api}

\subsection{What it does}

GENESYS is the model-based systems engineering package in which the client's
model lives. Its Client API is a set of .NET assemblies, namely
\code{Vitech.Genesys.Client} and \code{Vitech.Genesys.Common}, that let an
ordinary .NET program read and write the same repository the application works
on. The vendor's documentation describes it as giving external developers access
to the same interfaces the GENESYS application itself uses.

Through it a program can authenticate against a repository, open a project, read
the project schema at run time, and create, read, update and delete entities and
the relationships between them.

\subsection{Why it is useful}

It removes every intermediate step. Functions, flows, states and transitions can
be created from code and read back without exporting a file, parsing an
interchange format, or driving the user interface. The tool can turn C\# into
model content and model content into C\# while treating the repository as a
live object graph rather than as a document.

It also reads the schema rather than assuming it. Entity classes and attributes
are defined per project in GENESYS, so a tool that hard-codes class names breaks
on a customised schema. Reading the schema at run time avoids that.

Three specific uses in this project:

\begin{itemize}[nosep]
	\item Reading a state machine with its states, transitions, triggers and
	      guard conditions, which is the input to the state-to-activity translation.
	\item Reading function entities with their parameters and relationships,
	      which is the input to code generation.
	\item Creating the entities and relationships that an extracted control flow
	      graph maps onto, which is the output of code analysis.
\end{itemize}

Worked examples ship with the product, under \texttt{Samples\textbackslash API
Samples} in the installation directory. One of them,
\texttt{StructureTraversalSample}, walks a function's behaviour diagram using a
visitor over the construct types, which is directly the traversal the generation
half of this tool performs. It is the basis of the current
\code{Pentateuch.Genesys} test.

Documentation is the \emph{API Getting Started Guide} published by the vendor,
together with the full reference supplied as a compiled help file with the
GENESYS documentation package. The reference is the authoritative source for
whether a given operation exists, since the guide covers only the common
cases.

\subsection{Why it works for this project}

It is the only interface that is documented, supported, programmatic, and
complete. Every alternative in Chapter~\ref{ch:rejected} fails at least one of
those.

Two conditions come with it. The assemblies target .NET Framework 4.8, so the
code that calls them must too, which is the constraint Chapter~\ref{ch:runtime}
is a response to. And a licensed local installation must be present and running,
along with the Sentinel licensing components, which has implications for
handover.

\subsection{The REST API, and why it is not available}

GENESYS also exposes a REST API, added in the 2021 R2 release, which complements
the Client API. It is addressed over HTTP with token authentication, and third
party products integrate with GENESYS through it.

It is not available to this project. Enabling it requires server-side
configuration through the GENESYS administrative tools that the team does not
have rights to perform, and the client confirmed that access would not be
provided. 

This is worth mentioning due to the advantages it would've otherwise had.
It is transport-based rather than assembly-based, so it imposes no framework
constraint, and the whole tool could then have been a single modern .NET program
with no split and no shared contract library. The architecture in
Chapter~\ref{ch:runtime} exists in large part because this option is closed.

The practical effect is that the Client API is not one option among several. It
is the only programmatic route into the repository that the team can use.

\section{Roslyn}
\label{sec:roslyn}

\subsection{What it does}

Roslyn is the .NET compiler platform, and it exposes the compiler's inner workings
as an API. It parses source into a syntax tree that preserves every detail
including comments and whitespace, and it produces a semantic model that resolves
what the code actually means: which method a call resolves to, what type a
variable holds, and what paths control flow can take through a method.

It also works in the other direction. It can construct and edit syntax trees,
format the result, and compile it to check for errors. Visual Studio's own
navigation, refactoring and code suggestions are built on it, which is a
reasonable indication of its completeness and reliability.

\subsection{Why it is useful}

It covers roughly half the translation by itself, and both directions of it.

For model to code, the generation half can hand Roslyn a structure and have it
emit correctly formed, correctly formatted source, rather than assembling text
and hoping it parses. The compile check then answers whether the generated code
builds, which is a requirement in its own right.

For code to model, the semantic model is what makes the analysis possible at all.
A syntax tree alone would show that a call to \texttt{Process} exists; the
semantic model says which \texttt{Process}, in which type, in which assembly.
That resolution is what lets a call become a relationship between two entities
rather than a string. Roslyn also produces control flow graphs directly, which is
the structure the extracted activity is built from.

For keeping the two in step, the editing capability matters. Because Roslyn works
on a resolved view of a solution rather than on text, it can recognise that a
symbol was renamed rather than deleted and recreated, which is what a push and
pull reconciliation needs.

\subsection{Why it works for this project, and its caveats}

It is the right tool, and there is no serious alternative for analysing C\# with
semantic resolution. Four caveats attach to it.

\paragraph{Language coverage.}
Roslyn covers C\# and Visual Basic, which are the languages GENESYS supports for
its own scripting. If the project needs to translate C as well, Roslyn does not cover it 
and a second analysis front end would be required.

\paragraph{The in-process version conflict.}
GENESYS loads its own Roslyn to compile scripts. A tool running inside the
GENESYS process would be bound to that version, which is several years old, or
would break the host's scripting. This is the reason for the architecture in
Chapter~\ref{ch:runtime} rather than a reason against Roslyn itself.

\paragraph{The benefit of a current version.}
Running Roslyn outside GENESYS allows a current release, which can open whole
solutions rather than individual files, is substantially faster, and understands
recent C\# syntax while still analysing older code correctly. None of this would
be available in process.

\paragraph{Project resolution.}
Because Roslyn is a compiler rather than a parser, it wants a project that
resolves. Missing references or a solution that does not build produce error
types rather than resolved symbols, and the analysis degrades quietly unless
those are detected and reported.

\subsection{Packages used}

\begin{table}[htbp]
	\centering
	\caption{Roslyn packages}\label{tab:roslyn-packages}
	\begin{tabular}{L{0.34\textwidth}L{0.58\textwidth}}
\toprule
	\textbf{Package} & \textbf{Purpose} \\
	\midrule
\code{Microsoft.CodeAnalysis.CSharp} & Syntax trees, the
	semantic model and control flow analysis for C\#. The core of both translation
	directions. \\
	\addlinespace
	\code{Microsoft.CodeAnalysis.CSharp.Workspaces} & Project
	and document abstractions above a single file, including the editing and
	formatting services used when generating code. \\
	\addlinespace
	\code{Microsoft.CodeAnalysis.Workspaces.MSBuild} & Opens a
	solution or project file and resolves its references, which is what allows
	analysis across a whole codebase rather than file by file. This package is the
	main reason the analysis half targets a modern runtime. \\
	\bottomrule
	\end{tabular}
\end{table}

Versions should be pinned rather than floating, since Roslyn's behaviour changes
between releases and an analysis result that differs between developers would be
difficult to diagnose.

\section{Version Control and Project Management}

\subsection{Git and GitHub}

Git is used for version control and GitHub for hosting, issue tracking.

Three aspects are worth recording because they affect the design rather than 
the workflow. The code repository the tool reads and writes to is expected to be
under Git, which gives the users version control for the code without explicitly 
implementing it. The tool performs no version control operations itself, which 
keeps authentication, merge conflicts and branch management outside its scope. 
And reading from a commit rather than the working directory gives a stable snapshot 
rather than half-saved files.

\subsection{Branch protection}

A repository ruleset protects \texttt{main}. Direct pushes to it are blocked, so
every change arrives through a pull request, and a pull request requires an
approving review before it can merge.

\subsection{Task tracking}

GitHub Projects holds the backlog. Work is tracked there rather than in the issue
list alone, so that status is visible across the workstreams without
reading every issue.

\section{Documentation and Diagramming}

\subsection{\LaTeX}

Project documentation is written in \LaTeX, compiled on Overleaf. It gives
consistent formatting across documents, a structure that three people can edit in
parallel without merge conflicts, and cross-referencing that stays correct as
sections move.

\subsection{draw.io}

Diagrams are drawn in draw.io. It stores diagrams as XML, which means they diff
and merge in version control rather than being opaque binaries, and it exports to
SVG and PDF for inclusion in the documents. The editable sources are kept in the
repository beside the exported figures.

\subsection{OneNote}

Working notes are kept in OneNote: meeting notes with the client, findings while
investigating the GENESYS interfaces, and the rough material that later becomes a
section of one of these documents.

\section{Supporting Development Tools}

\begin{table}[htbp]
	\centering
	\caption{Supporting tools}\label{tab:supporting}
	\begin{tabular}{L{0.24\textwidth}L{0.68\textwidth}}
\toprule
	\textbf{Tool} & \textbf{Use} \\
	\midrule
Visual Studio 2026 & Primary development environment, with the .NET desktop
	development workload, which supplies both .NET Framework 4.8 and the .NET 10
	SDK. Required in practice for the .NET Framework half, where command line
	tooling resolves framework references less reliably. \\
	\addlinespace
	.NET SDK & Command line build and run across all three projects. \\
	\addlinespace
	NuGet & Package management, principally for the Roslyn packages, which restore
	automatically on first build. \\
	\bottomrule
	\end{tabular}
\end{table}
